% Generated draft: requires mathematical and source review.
\documentclass[11pt]{article}
\usepackage[margin=25mm]{geometry}
\usepackage{fontspec}
\setmainfont{Latin Modern Roman}
\usepackage{amsmath,amssymb,mathtools,unicode-math}
\setmathfont{Latin Modern Math}
\usepackage{enumitem,needspace,xurl,hyperref}
\setlength{\parindent}{0pt}
\setlength{\parskip}{5pt}
\setlength{\emergencystretch}{3em}
\newcommand{\N}{\mathbb{N}}
\newcommand{\Z}{\mathbb{Z}}
\newcommand{\Q}{\mathbb{Q}}
\newcommand{\R}{\mathbb{R}}
\newcommand{\C}{\mathbb{C}}
\newcommand{\F}{\mathbb{F}}
\newcommand{\A}{\mathbb{A}}
\newcommand{\PP}{\mathbb{P}}
\newcommand{\E}{\mathbb{E}}
\newcommand{\eps}{\varepsilon}
\newcommand{\dd}{\,\mathrm{d}}
\begin{document}
\section*{UnsolvedMath ID: 30000995}
\textbf{Source problem code:} catalogue entry
\subsection{Definitions and mathematical statement}
The cleaned source asks: given a finite set \(P\) of points in \(\R^2\) or \(\R^3\), what is the maximum number of these points lying on the graph of a \(1\)-Lipschitz function?

A precise conventional model, conditional on using the last coordinate as height and Euclidean norms, is the following. For \(d\in\{2,3\}\), write \(\R^d=\R^{d-1}\times\R\). A function \(f:\R^{d-1}\to\R\) is \(1\)-Lipschitz if
\[
|f(x)-f(y)|\leq\|x-y\|_2\quad\hbox{for all }x,y\in\R^{d-1}.
\]
Determine, from the given finite set \(P\subset\R^d\), the quantity
\[
M(P)=\max_{\operatorname{Lip}(f)\leq1}
\bigl|P\cap\{(x,f(x)):x\in\R^{d-1}\}\bigr|.
\]
The extraction does not specify the height direction, norm, domain of \(f\), or whether an algorithm, deterministic bound, or random-model asymptotic is intended. The displayed model therefore explains one reading rather than fixing these missing choices as source hypotheses. The original also says ``finite distribution or set''; a weighted or random distribution would need its own objective and probability law. No such distribution is supplied.
\subsection{Short English statement}
For a given finite planar or spatial point set, maximize the number of points on a \(1\)-Lipschitz graph. The natural finite-set optimization can be written down once the graph direction and metric are fixed, but the source's distributional and boarding-policy context is incomplete.
\subsection{Sources}
\begin{itemize}
\item Local UnsolvedMath registry, permanent ID \(30000995\), source code \texttt{catalogue entry}, ``Point Incidence on One-Lipschitz Graphs.'' Registry curation: Ulam AI / UnsolvedMath Contributors, CC BY 4.0; underlying sources retain their supplied provenance.
\item \emph{Calculus of Variations} (2009), Oberwolfach Reports, pp.\ 1707--1770, DOI \url{https://doi.org/10.4171/owr/2008/31}.
\item Supplied literature triage, \texttt{literature\_checked\_at}: 2026-08-21: labeled open, but says the extracted passage stops before specifying the boarding-policy distribution or optimization problem and is incomplete rather than a verifiable solved claim. This is dated registry curation and requires independent review; its boarding-policy wording does not match the complete finite-set optimization suggested by the cleaned statement.
\end{itemize}
\end{document}
